\SixthArticle{24}
Da der Grenzwert der Funktion $F(\alpha,\beta,\gamma+k,1)$, wenn $k$ ins Unendliche wächst, offenbar die Einheit ist, geht Gleichung 39 in folgende über:
\[
\begin{array}{@{}ll@{}}
[48]&\displaystyle F(\alpha,\beta,\gamma,1)=
\frac{\Pi(\gamma-1).\Pi(\gamma-\alpha-\beta-1)}
{\Pi(\gamma-\alpha-1).\Pi(\gamma-\beta-1)}
\end{array}
\]
Diese Formel gibt die vollständige Lösung der Frage, die den Gegenstand dieses Abschnitts bildet. Daraus folgen unmittelbar die eleganten Gleichungen:
\SourceLeaf{157}
\[
\begin{array}{@{}ll@{}}
[49]&\displaystyle F(\alpha,\beta,\gamma,1)=F(-\alpha,-\beta,\gamma-\alpha-\beta,1)\\[3pt]
[50]&\displaystyle F(\alpha,\beta,\gamma,1).F(-\alpha,\beta,\gamma-\alpha,1)=1\\[3pt]
[51]&\displaystyle F(\alpha,\beta,\gamma,1).F(\alpha,-\beta,\gamma-\beta,1)=1
\end{array}
\]
wobei in der ersten $\gamma$, in der zweiten $\gamma-\beta$ und in der dritten $\gamma-\alpha$ eine positive Größe sein muss.

\SixthArticle{25}
Wenden wir Formel 48 auf einige der Gleichungen des Artikels 5 an. Setzt man $t=90^0=\frac12\pi$, so wird Formel XIII zu $\frac12\pi=F(\frac12,\frac12,\frac32,1)$ oder ist gleichbedeutend mit der bekannten Gleichung
\[
\frac12\pi=1+\frac{1.1}{2.3}+\frac{1.1.3}{2.4.5}
+\frac{1.1.3.5}{2.4.6.7}+\mathrm{usw.}
\]
Da nun nach Formel 48 $F(\frac12,\frac12,\frac32,1)=\frac{\Pi\frac12.\Pi(-\frac12)}{\Pi0.\Pi0}$ und außerdem $\Pi0=1$, $\Pi\frac12=\frac12\Pi(-\frac12)$ gilt, ergibt sich $\pi=(\Pi(-\frac12))^2$ oder
\[
\begin{array}{@{}ll@{}}
[52]&\displaystyle\Pi(-\tfrac12)=\BareRoot{\pi}\\[3pt]
[53]&\displaystyle\Pi\tfrac12=\tfrac12\BareRoot{\pi}
\end{array}
\]
Formel XVI des Artikels 5, die gleichbedeutend mit der bekannten Gleichung
\[
\sin nt=n\sin t-\frac{n(nn-1)}{2.3}\sin t^3
+\frac{n(nn-1)(nn-9)}{2.3.4.5}\sin t^5-\mathrm{usw.}
\]
ist und allgemein für jeden beliebigen Wert von $n$ gilt, sofern $t$ die Grenzen $-90^0$ und $+90^0$ nicht überschreitet, ergibt für $t=\frac12\pi$
\[
\sin\frac{n\pi}{2}=
\frac{n\Pi\frac12.\Pi(-\frac12)}{\Pi(-\frac12n).\Pi\frac12n}
\]
woraus die elegante Formel folgt:
\[
\Pi\frac12n.\Pi(-\frac12n)=\frac{\frac12n\pi}{\sin\frac12n\pi},
\quad\hbox{oder mit }n=2z
\]
\[
\begin{array}{@{}ll@{}}
[54]&\displaystyle\Pi(-z).\Pi(+z)=\frac{z\pi}{\sin z\pi}\\[3pt]
[55]&\displaystyle\Pi(-z).\Pi(z-1)=\frac{\pi}{\sin z\pi}
\end{array}
\]
und außerdem, wenn $z+\frac12$ für $z$ geschrieben wird,
\[
\begin{array}{@{}ll@{}}
[56]&\displaystyle\Pi(-\tfrac12+z).\Pi(-\tfrac12-z)=\frac{\pi}{\cos z\pi}
\end{array}
\]
\SourceLeaf{158}
Aus der Verbindung der Formel 54 mit der Definition der Funktion $\Pi$ folgt, dass $\frac{z\pi}{\sin z\pi}$ der Grenzwert des unendlichen Produkts
\[
\frac{(1.2.3.4\,.\,.\,.\,.k)^2}
{(1-zz)(4-zz)(9-zz)\,.\,.\,.\,.(kk-zz)}
\]
ist, wenn $k$ ins Unendliche wächst; daher ist
\[
\sin z\pi=z\pi(1-zz)(1-\frac{zz}{4})(1-\frac{zz}{9})\quad
\mathrm{usw.\ ins\ Unendliche}
\]
In gleicher Weise ergibt sich aus 56
\[
\cos z\pi=(1-4zz)(1-\frac{4zz}{9})(1-\frac{4zz}{25})\quad
\mathrm{usw.\ ins\ Unendliche}
\]
Dies sind die sehr bekannten Formeln, die von den Analysten gewöhnlich mit völlig anderen Methoden hergeleitet werden.

\SixthArticle{26}
Bezeichnet $n$ eine ganze Zahl, so ergibt sich der Wert des Ausdrucks
\[
\frac{n^{nz}\Pi(k,z).\Pi(k,z-\frac1n).\Pi(k,z-\frac2n)
\,.\,.\,.\,.\Pi(k,z-\frac{n-1}{n})}{\Pi(nk,nz)}
\]
nach gehörigem Zusammenfassen zu
\[
=\frac{(1.2.3\,.\,.\,.\,.k)^n n^{nk}}
{1.2.3\,.\,.\,.\,.nk.k^{\frac12(n-1)}}
\]
und ist daher von $z$ unabhängig, bleibt also gleich bei jeder Wertzuweisung an $z$. Da $\Pi(k,0)=\Pi(nk,0)=1$ ist, lässt er sich somit als folgendes Produkt darstellen:
\[
\Pi(k,-\frac1n).\Pi(k,-\frac2n).\Pi(k,-\frac3n)
\,.\,.\,.\,.\Pi(k,-\frac{n-1}{n})
\]
Lassen wir also $k$ ins Unendliche wachsen, so erhalten wir
\[
\begin{gathered}
\frac{n^{nz}\Pi z.\Pi(z-\frac1n).\Pi(z-\frac2n)
\,.\,.\,.\,.\Pi(z-\frac{n-1}{n})}{\Pi nz}\\[4pt]
=\Pi(-\frac1n).\Pi(-\frac2n).\Pi(-\frac3n)
\,.\,.\,.\,.\Pi(-\frac{n-1}{n})
\end{gathered}
\]
Multipliziert man das Produkt auf der rechten Seite mit sich selbst bei umgekehrter Reihenfolge der Faktoren, so ergibt sich nach Formel 55
\[
\frac{\pi}{\sin\frac1n\pi}.\frac{\pi}{\sin\frac2n\pi}.
\frac{\pi}{\sin\frac3n\pi}\,.\,.\,.\,.
\frac{\pi}{\sin\frac{n-1}{n}\pi}=\frac{(2\pi)^{n-1}}{n}
\]
Daher haben wir den eleganten Satz
\SourceLeaf{159}
\[
\begin{aligned}
{}\text{[57]}\quad &
\frac{n^{nz}\Pi z.\Pi(z-\frac1n).\Pi(z-\frac2n).\,.\,.\,.
\Pi(z-\frac{n-1}{n})}{\Pi nz}
=\frac{(2\pi)^{\frac12(n-1)}}{\BareRoot n}
\end{aligned}
\]

